\section{Budgeted Data Selection as Bilevel Optimization}
\label{sec:formulation}

Let $\Dtr=\{z_i\}_{i=1}^{n}$ be candidate training examples and let
$\Dval$ be a clean holdout set. We use
\begin{align}
g(\vomega,\vtheta)
    &:= \frac{1}{n}\sum_{i=1}^{n}\omega_i\ell(\vtheta;z_i),
    &
f(\vtheta)
    &:= \frac{1}{|\Dval|}\sum_{z\in\Dval}\ell(\vtheta;z),
\label{eq:fg}
\end{align}
and define
\begin{equation}
\min_{\vomega\in\Wbudget} h(\vomega)
    := f(\thetahat(\vomega))
\quad\text{subject to}\quad
\thetahat(\vomega)\in
    \argmin_{\vtheta}g(\vomega,\vtheta),
\label{eq:bilevel-selection}
\end{equation}
where
\begin{equation}
\Wbudget:=\left\{\vomega\in[0,1]^n:
\sum_{i=1}^{n}\omega_i=B\right\}.
\label{eq:capped-simplex}
\end{equation}
The continuous vector $\vomega$ is the outer optimization state. The target
model is trained on the discrete set
$S(\vomega)=\operatorname{TopB}(\vomega)$ only when a training decision is
required. The binary mask is not substituted back for $\vomega$ in the next
outer update.

The budget-feasible uniform point is
\begin{equation}
\vomega^{\mathrm{unif}}=\frac{B}{n}\bm{1}.
\label{eq:uniform-budget}
\end{equation}
It is the appropriate reference for comparing one-shot and iterative
selection. The all-ones vector has a useful attribution interpretation---all
examples are fully present---but lies outside $\Wbudget$ whenever $B<n$.
Conflating these two reference points creates an artificial first projection
that can dominate the score update.

Equation~\ref{eq:bilevel-selection} is our optimization target, not a claim
that the nonconvex neural-network problem satisfies the assumptions used in
all bilevel convergence analyses. In particular, the F2SA theory assumes an
unconstrained outer variable and a strongly convex lower problem. We use its
value-function construction to derive an estimator, and treat projection onto
the set in Eq.~(\ref{eq:capped-simplex}) as an explicit algorithmic adaptation whose behavior
must be tested empirically.

F2SA's MNIST hyper-cleaning experiment instead optimizes unconstrained logits
$\va\in\mathbb{R}^n$ and uses sigmoid weights $\omega_i=\sigma(a_i)$; it does
not impose a fixed cardinality budget. To separate parameterization from the
estimator, our primary method optimizes $\vomega$ directly with capped-simplex
projection. As an ablation, we use
\begin{equation}
\omega_i=\sigma(a_i+c(\va)),
\qquad \sum_i\sigma(a_i+c(\va))=B,
\end{equation}
where the scalar shift $c(\va)$ is obtained by one-dimensional root finding.
Learning rates are tuned separately because direct weights and logits have
different geometries.
